\scanpage{191}
% l. 01
( this \unsure{works} \unsure{without} a finiteness hypothesis on
% l. 02
$\Gamma$, the universal case being the one where $\Gamma$
% l. 03
is the Teichmüller \unsure{group} of $\pi'$ \unsure{fixing} $i$
% l. 04
i.e. $\operatorname{Autext}_{\mathrm{lac}}(\pi';i)\simeq\operatorname{Aut}_{\mathrm{lac}}(\pi)$ ). I
% l. 05
say that if $\Gamma$ operates \uline{faithfully} \uline{on $\pi$}
% l. 06
i.e. if its outer operation \uline{on}
% l. 07
\uline{$\pi'$} is faithful, and if $\Gamma$ is \uline{finite},
% l. 08
Then $\Gamma$ is necessarily cyclic \inserted{(case $\Gamma=\Gamma^0$)}
% l. 09
or dihedral ($\Gamma\ne\Gamma^0$), and that if in
% l. 10
addition $\Gamma=\Gamma^0$, then \deleted{\unsure{$\mathcal E_i'$ is $\simeq$}} \inserted{\unsure{$\mathcal E_i\simeq\mathbb Z$|$\mathcal E_i'\simeq\mathbb Z$} i.e.} $\exists!$ isomorphism
% l. 11
\unsure{$T\xrightarrow{\sim}\mathcal E_i$|$T\xrightarrow{\sim}\mathcal E_i'$} such that \unsure{$T\xrightarrow{k_i'}\mathcal E_i'$} identifies with
% l. 12
$n\operatorname{id}_T$ (so that we find \unsure{$\Gamma\simeq\mathcal E_i'/n\mathcal E_i'$}
% l. 13
$\simeq T\otimes_{\mathbb Z}\mathbb Z/n\mathbb Z$, if $n=\operatorname{card}\Gamma$ ...

% l. 14
Changing notation, this amounts to the
% l. 15
\uline{Theore}\textsuperscript{\unsure{m}} Let $\Gamma$ be a finite group operating
% l. 16
faithfully on an \uline{\unsure{outer|outer\textsuperscript{ly}}}
% l. 17
group with loops $\pi$ \inserted{of type \deleted{\illegible{1}} $(g,\nu+1)$ ($\nu\geq0$)}, leaving an $i\in I(\pi)$ fixed,
% l. 18
Then $\Gamma$ is cyclic (if $\Gamma=\Gamma^0$) or
% l. 19
dihedral (if $\Gamma\ne\Gamma^0$), and in the extension
% marge, en bas à gauche
\margin{\unsure{to the} same there is \uline{equivalence}, if in the theorem, we assume the type $(g,\nu)$ anabelian --- otherwise the theorem is a little more general}
\mspage{190}


\scanpage{192}
\mspage{96}
% l. 01
corresponding \deleted{\illegible{2}} \unsure{$\mathcal E$} of $\Gamma$ by $\pi$,
% l. 02
If \unsure{$\mathcal E_i$} is the normalizer of an $L_i$ in
% l. 03
\unsure{$\mathcal E$} [ \inserted{so that we have} \deleted{an} \unsure{$\mathcal E_i\to\Gamma$} epimorphic i.e. an exact sequence
% l. 04
\unsure{$1\to L_i\to\mathcal E_i\to\Gamma\to1$} ]. the inverse image \unsure{$\mathcal E_i^0$} of
% l. 05
$\Gamma^0$ is $\simeq\mathbb Z$ so (if $n=\operatorname{card}\Gamma^0$)\inserted{\unsure{$=[\mathcal E_i^0:L_i]$}}
% l. 06
\deleted{\unsure{$n\operatorname{id}$}} $x\to x^n$ is an isomorphism \deleted{and} \unsure{$\mathcal E_i^0\simeq L_i$}, which
% l. 07
taking $k_i$ into account gives an
% l. 08
isomorphism \unsure{$\mathcal E_i^0\simeq T$}, whose composite with
% l. 09
\deleted{\unsure{$\mathcal E_i$}} \unsure{$T\xrightarrow[\sim]{k_i}L_i\subset\mathcal E_i^0$} is $n\operatorname{id}_T$, so
% l. 10
that we have a canonical isomorphism
% l. 11
\[\unsure{\Gamma^0\simeq\mathcal E_i^0/L_i\simeq T/nT}\]
% l. 12
(obviously independent of the choice of $L_i$ ...)

% l. 13
\uline{Proof}. Let us consider the extension
% l. 14
\unsure{$\mathcal E_i^0$} of $\Gamma^0$ by $L_i\simeq T$ ($\overset{\text{\unsure{noncanonical}}}{\simeq}\mathbb Z$),
% l. 15
since \unsure{$\mathcal E_i^0$|$\Gamma^0$} operates trivial\textsuperscript{ly} on $T$ \inserted{torsion-free}, this
% l. 16
extension (by the cohomology exact sequence
% l. 17
associated with $0\to T\to T\otimes_{\mathbb Z}\mathbb Q\to T\otimes_{\mathbb Z}\mathbb Q/\mathbb Z\to0$)
% l. 18
is canonical\textsuperscript{ly} isomorphic to \deleted{\unsure{an}} the extension
% l. 19
defined by a homomorphism $\Gamma^0\to T_n=T/nT$,
% l. 20
as inverse image of the extension
% l. 21
$0\to T\xrightarrow{n}T\to T_n\to0$ of $T_n$ by $T$. I say that
\mspage{191}


\scanpage{193}
% l. 01
this homomorphism is an isomorphism (from which \unsure{would follow|will follow}
% l. 02
the other assertions), or what amounts
% l. 03
to the same since $\Gamma^0$ and $T_n$ are both
% l. 04
of order $n$, that it is a monomorphism.
% l. 05
Replacing $\Gamma^0$ by the kernel of $\Gamma^0\to\mathbb Z/n\mathbb Z$,
% l. 06
we may assume that the homomorphism \inserted{in question} is zero, and we must
% l. 07
prove that this implies that the action
% l. 08
of $\Gamma^0$ is trivial (which, since by hypothesis
% l. 09
the action of $\Gamma$ is faithful, implies $\Gamma^0=\{1\}$,
% l. 10
OK.) So \unsure{theorem} is reduced to the

% l. 11
\uline{fundamental lemma}. Every \inserted{direct} outer automorphism
% l. 12
$u$ of finite order \inserted{\unsure{$n$}} of a group with
% l. 13
loops $\pi$, which fixes a loop class
% l. 14
$i$ and is such that the extension of $\mathbb Z/n\mathbb Z$
% l. 15
by $L_i\in i$ defined by $u$ is trivial,
% l. 16
is trivial.

% l. 17
The hypothesis means that $u$ lifts
% l. 18
to an automorphism $u_0$ of $\pi$, which
% l. 19
normalizes $L_i$, and which is also of order \unsure{$n$},
% l. 20
(or of finite order, this amounts to the \unsure{same|$\Rightarrow$} taking into account
% l. 21
that \deleted{\unsure{the elements of}} $T$ has no elements of finite order)
\mspage{192}


\scanpage{194}
\mspage{97}
% l. 01
then $u_0$ is trivial, i.e. this is equivalent to the
% l. 02
\uline{Corollary}. Every \inserted{direct} automorphism of finite order \inserted{$u_0$} of
% l. 03
a group with loops $\pi$, \inserted{which is direct and} which normalizes
% l. 04
a loop subgroup $L_i$ (i.e. which
% l. 05
centralizes $L_i$), is reduced to the identity.
% marge supérieure gauche
\margin{In the lemma or its \unsure{corollary}}
\margin{\uline{N.B.} (the case $g=0$, $\nu+1=1$ or $\nu+1=2$ is trivial, the case $\nu+1=3$ ($(g,\nu)=(0,2)$ abelian!) is not trivial \unsure{on the other hand}, nor the case $g=1$, $\nu+1=1$ i.e. $(g,\nu)=(1,0)$ abelian. Yet the result should be valid also in \unsure{this|these} cases.)}

% l. 06
To prove it, we are reduced \inserted{at once} to the
% l. 07
case where $u_0$ is \deleted{of prime order $p$} \inserted{such that $u_0^p=\operatorname{id}$, where}
% l. 08
$p$ is a prime \unsure{number} --- i.e. $u_0$ corres-
% l. 09
ponds to an operation in the strict sense
% l. 10
(not on\textsuperscript{ly} outer) of $\mathbb Z/p\mathbb Z$ \deleted{\unsure{by}} \unsure{by automorphisms}
% l. 11
of $\pi$. But (whether or not we can
% l. 12
find such a $p$) let us consider the case where
% l. 13
we know that the \uline{outer} operation of
% l. 14
$\Gamma=\mathbb Z/n\mathbb Z$ [$n=$ order of $u$] on $\pi$ is realized
% l. 15
geometrical\textsuperscript{ly}, by a (faithful) operation
% l. 16
of $\Gamma=\mathbb Z/n\mathbb Z$ on a $U$ of type $(g,\nu+1)$.
% l. 17
\deleted{\unsure{$\pi$-conjugacy classes}, \illegible{1}}
% marge inférieure gauche
\margin{$U=X\setminus S$, $S\simeq I$, $X$ compact of genus $g$ \unsure{the puncture|the part} $s_i$ of $S$ corresponding to a fixed point of $\Gamma$ operating on $U$.}

% l. 18
\deleted{We know that the (\unsure{liftings}) of the outer automorphisms
% l. 19
to \deleted{\unsure{automorphisms}} \inserted{a genuine operation},
% l. 20
correspond to the fixed points \illegible{1} of $\Gamma$
% l. 21
in $U$; \deleted{and the extension of} \inserted{let $a$ be the fixed point of $\Gamma$}
% l. 22
corresponding to the given operation, \unsure{so}
% l. 23
that we have \deleted{\unsure{a}} canonical isomorphism \unsure{compatible with $\Gamma$}
% l. 24
$\pi=\pi_1(U,a)$.}
\mspage{193}


\scanpage{195}
% l. 01
We express the hypothesis of the \inserted{\unsure{``genuine''}} operation of $\Gamma$ on
% l. 02
$\pi$, \deleted{\illegible{1}} \unsure{centralizing} an $L_i$, by saying that
% l. 03
the outer operation gives a trivial extension
% l. 04
of $\Gamma$ by $L_i\simeq T$. But we know
% l. 05
moreover in the geometric case \inserted{and faithful operation} that it
% l. 06
is \uline{not} trivial !

% l. 07
It would therefore suffice, \deleted{\illegible{1}} to prove the
% l. 08
fundamental lemma, to know that every
% l. 09
outer action \inserted{[faithful]} of a cyclic group
% l. 10
on a group with loops of type \unsure{$(g,\nu+1)$}
% l. 11
is realizable, and it even suffices to know this
% l. 12
for a cyclic group of
% l. 13
prime order. \inserted{\unsure{is}} \unsure{Now} unless I am mistaken this result
% l. 14
is \unsure{well} known (even for \unsure{solv-}
% l. 15
able groups ...). --- so that the funda-
% l. 16
mental lemma seems proved. I have only
% l. 17
a doubt whether it is proved in the
% l. 18
general case of a $(g,\nu)$, \inserted{\deleted{\illegible{1} $g\geq2$}} or only
% l. 19
for \inserted{$g\geq2$,} $\nu=0$. But \deleted{in truth, \illegible{2}, \inserted{this suffices}}
% l. 20
\deleted{\illegible{2} \inserted{for $g\geq2$}}
% marge inférieure gauche
\margin{as a theorem \unsure{of existence} \unsure{of points} fixed by a group \illegible{1} \unsure{operating} \unsure{of} Teichmüller}

% l. 21
\deleted{\uline{Lemma} Let \illegible{1} be an outer operation}
\mspage{194}